\documentclass[11pt]{article}
\usepackage[T1]{fontenc}
\usepackage[utf8]{inputenc}
\usepackage{lmodern,amsmath,amssymb,amsthm,mathtools}
\usepackage[margin=1in]{geometry}
\usepackage{booktabs,microtype}
\usepackage[numbers,sort&compress]{natbib}
\usepackage{url}
\usepackage[hidelinks]{hyperref}
% Certificate files are attached by the release build driver.
\hypersetup{pdftitle={Binomial log-concavity of the lower-mode forcing term in a rank-two Chern recurrence},
pdfauthor={Xiaojian Zeng},pdfsubject={Auxiliary log-concavity theorem with an exact polynomial certificate}}
\newtheorem{theorem}{Theorem}
\newtheorem{lemma}[theorem]{Lemma}
\newcommand{\LC}{\operatorname{LC}}
\newcommand{\supp}{\operatorname{supp}}
\DeclareMathOperator{\sort}{sort}
\allowdisplaybreaks[1]
\title{Binomial log-concavity of the lower-mode forcing term\\
in a rank-two Chern recurrence}
\author{Xiaojian Zeng\\[3pt]
\small PBC School of Finance, Tsinghua University, 100190\\
\small \href{mailto:xiaojianzeng@foxmail.com}{xiaojianzeng@foxmail.com}}
\date{14 September 2026}
\begin{document}
\maketitle
\begin{abstract}
The rank-two binomial log-concavity conjecture of B\'erczi and Feh\'er
concerns the Newton coefficients of Schur coefficients of symmetric-power
Chern classes. We prove an auxiliary theorem for the sum of the four
lower-mode forcing terms in their recurrence for the second Schur mode.
For every $k\geq4$, this forcing polynomial has positive Newton
coefficients exactly on $[k-2,2k-1]$, and the coefficient sequence is
log-concave. The proof applies the prior permutation-defect and Abel
summation method, with a complete certificate of 41 polynomial prefixes.
Their nonnegative expansions are checked in exact arithmetic, and the
data and reconstruction program are included as PDF attachments.
The original conjecture and its second-mode restriction remain unresolved
by this result; we identify a sufficient mixed inequality for that restriction.
\end{abstract}

\section{The auxiliary theorem and its scope}

For integers $d,k\geq0$, define
\begin{equation}\label{eq:chern}
c_k(2,d)=[z^k]\prod_{i=0}^{d}\bigl(1+z(ix_1+(d-i)x_2)\bigr).
\end{equation}
For $0\leq j\leq\lfloor k/2\rfloor$, write
\begin{equation}\label{eq:schur}
c_k(2,d)=\sum_j A_{k,j}(d)s_{(k-j,j)}(x_1,x_2),\qquad
s_{(k-j,j)}=\sum_{m=j}^{k-j}x_1^{k-m}x_2^m.
\end{equation}
The functions $A_{k,j}$ are rational polynomials in $d$
\cite[Theorem 2.4 and Section 3]{BercziFeher2026}. Their Newton expansions are
\begin{equation}\label{eq:newton}
A_{k,j}(x)=\sum_{r\geq0} B_{k,j,r}\binom{x}{r},\qquad
B_{k,j,r}=\Delta^r A_{k,j}(0),\qquad
\Delta f(x)=f(x+1)-f(x).
\end{equation}
Here $\binom{x}{r}=x(x-1)\cdots(x-r+1)/r!$, with $\binom{x}{0}=1$.
All coefficient sequences below are extended by zero at negative indices.
Polynomiality also follows directly from the Chern-root power sums:
their coefficients are sums of polynomials in $i,d$, hence polynomials
in $d$; Newton's identities and the triangular change in
\eqref{eq:schur} preserve polynomiality. Their degrees are at most $2k$.

B\'erczi--Feh\'er's Conjecture 3.3 asks for interval support and
$B_{k,j,r}^2\geq B_{k,j,r-1}B_{k,j,r+1}$ for every admissible $k,j,r$.
Their formal Theorems 3.7 and 5.4 prove the cases $j=0,1$, and
Theorem 5.2 proves positivity for $j=2$ \cite{BercziFeher2026}.
The theorem here concerns a different, explicitly defined polynomial.

For $k\geq4$, set
\begin{equation}\label{eq:residual}
\begin{aligned}
R_k(x)={}&(x-k+2)A_{k-1,1}(x)
 +\frac{(x-k+3)(x-k+2)}2 A_{k-2,0}(x)\\
&+(x+1)(x-k+3)A_{k-2,1}(x)\\
&+\frac{(x+1)(x-k+4)(x-k+3)}2 A_{k-3,0}(x),
\qquad R_k(x)=\sum_{r\geq0}\rho_{k,r}\binom{x}{r}.
\end{aligned}
\end{equation}

\begin{theorem}\label{thm:main}
For every integer $k\geq4$, the coefficients in \eqref{eq:residual}
are positive exactly for $k-2\leq r\leq2k-1$ and vanish otherwise.
For every integer $r\geq0$,
\begin{equation}\label{eq:main}
\rho_{k,r}^2\geq\rho_{k,r-1}\rho_{k,r+1}.
\end{equation}
\end{theorem}

The forcing polynomial \eqref{eq:residual} is already present in
\cite[equation (5.2)]{BercziFeher2026}:
\begin{equation}\label{eq:recurrence}
\Delta A_{k,2}(x)=(x+1)A_{k-1,2}(x)+R_k(x).
\end{equation}
At $k=4$ the out-of-range term $A_{3,2}$ means zero. This convention
follows by extracting $[u^2](1-u)$ from a palindromic cubic, and does not
change the admissible definitions in \eqref{eq:schur}.
For completeness, \eqref{eq:recurrence} follows by comparing powers of $T$
and then extracting $[u^2](1-u)$ in
\[
E_d(T,u):=\prod_{i=0}^d(T+i+(d-i)u)
=(T+d)E_{d-1}(T+u,u).
\]

The increment over the prior recurrence and positivity theorem is
log-concavity of the \emph{sum} of its four lower-mode terms. Log-concavity
does not in general survive addition: $(1,1,1,1)$ and $(1,2,4,8)$ are
positive log-concave sequences, whereas their sum $(2,3,5,9)$ is not.
Our certificate extends the 17-prefix calculation for $j=1$ in
\cite[Theorem 5.4]{BercziFeher2026} to this 41-prefix forcing family;
the recurrence and the general method are prior results.
Patel's reader contribution \cite{Patel2026} claims all-index binomial
positivity by a support induction. We use no unverified all-index
assertion from that contribution, and claim no such positivity result
as an increment here.

No equivalent theorem for the explicit $R_k$ was identified in the
inspected sources as of 14 September 2026. This is a bounded comparison,
not a claim of exhaustive priority verification. In comparing with
\cite{BercziFeher2026}, we use the formal $j=1$ statement and proof of
Theorem 5.4; the $j=2$ log-concavity label in its Figure 1 is inconsistent
with that statement. Neither the full original conjecture nor its $j=2$
restriction is proved here.

\section{Permutation coefficients and exact support}

Let $D(N,K)$ count permutations of $N$ letters with no fixed points and
$N-K$ cycles. Put $D(0,0)=1$ and set $D$ to zero outside its natural
range. Deleting a distinguished letter from a cycle of length at least
three, or deleting its transposition, gives
\begin{equation}\label{eq:defect}
D(N,K)=(N-1)\bigl(D(N-1,K-1)+D(N-2,K-1)\bigr).
\end{equation}
For $K\geq1$, its positive support in $N$ is exactly $[K+1,2K]$.
Indeed, $c=N-K$ nontrivial cycles require $c\geq1$ and $2c\leq N$.
Conversely, use $c-1$ transpositions and one cycle of length
$2K-N+2\geq2$.

Put $S_h(x)=A_{h,0}(x)$. At a nonnegative integer $d$,
$S_h(d)=e_h(0,1,\ldots,d)$ counts permutations of $d+1$ letters with
cycle defect $h$. Decomposition according to nonfixed support and
Pascal's identity give
\begin{equation}\label{eq:zeroth}
S_h(x)=\sum_{r\geq0}\bigl(D(r,h)+D(r+1,h)\bigr)\binom{x}{r}.
\end{equation}
An identity valid at all sufficiently large integer $d$ is a polynomial
identity in $x$. Similarly, the rising-factorial identity
\[
E_d(T,u)=(1-u)^{d+1}
\left(\frac{T+du}{1-u}\right)^{\overline{d+1}}
\]
and extraction of $[u](1-u)$ give
\begin{equation}\label{eq:first}
A_{n,1}(x)=x(x+2-n)S_{n-1}(x)-(n+1)S_n(x).
\end{equation}
These formulas and the defect basis are used in
\cite[Section 3]{BercziFeher2026}.

If $F(x)=\sum_r f_r\binom{x}{r}$, multiplication by $x-c$ gives
\begin{equation}\label{eq:linear}
\left[\binom{x}{r}\right](x-c)F(x)=(r-c)f_r+r f_{r-1},
\end{equation}
where the bracket denotes coefficient in the Newton basis. Substituting
\eqref{eq:zeroth} into \eqref{eq:first} and applying
\eqref{eq:defect} yields
\begin{equation}\label{eq:firstcoeff}
\begin{aligned}
B_{n,1,r}={}&r(r+2-n)D(r+1,n-1)
 +r(3r+2-3n)D(r,n-1)\\
&+(3r^2-3nr-2r+n+1)D(r-1,n-1)\\
&+(r-1)(r-n-1)D(r-2,n-1).
\end{aligned}
\end{equation}
Every multiplier is nonnegative whenever its $D$ factor is positive.
For the third multiplier, write $r=n+1+t$; it becomes
$2n+2+(3n+4)t+3t^2>0$ for $t\geq0$. The other multipliers are positive
on their active supports. For $n\geq2$, the first term covers
$n-1\leq r\leq2n-3$; the next three terms cover $r=2n-2,2n-1,2n$,
with multipliers $(2n-2)(3n-4)$, $6(n-1)^2$, and $(2n-1)(n-1)$.
Thus the strict support of $B_{n,1,r}$ is $[n-1,2n]$.
Equation \eqref{eq:zeroth} gives strict support $[h,2h]$ for $S_h$,
for every $h\geq1$.

If a nonnegative sequence is positive exactly on $[I,M]$ and $c\leq I$,
\eqref{eq:linear} gives strict support $[I,M+1]$ when $c<I$, and
$[I+1,M+1]$ when $c=I$. Applying this rule successively to the four
terms of \eqref{eq:residual} gives, in their displayed order,
\[
[k-1,2k-1],\quad [k-1,2k-2],\quad
[k-2,2k-2],\quad [k-2,2k-3].
\]
All four coefficient sequences are nonnegative. Their union proves the
strict-support assertion of Theorem~\ref{thm:main}, including $k=4$.

\section{The Abel certificate}

\begin{lemma}\label{lem:defectLC}
For fixed $K\geq1$, the sequence $D(N,K)$ is log-concave and, with zero
extension, satisfies
\[
D(p,K)D(q,K)\geq D(p-1,K)D(q+1,K)\qquad(p\leq q).
\]
\end{lemma}
\begin{proof}
Let $P_K(t)=\sum_N D(N,K)t^N$. Equation \eqref{eq:defect} gives
\[
P_K(t)=t^2\bigl((1+t)P_{K-1}(t)\bigr)',\qquad P_0(t)=1.
\]
Induction and Rolle's theorem show that $P_K$ has only real nonpositive
zeros. The usual coefficient inequalities for a real-rooted polynomial
with nonnegative coefficients imply log-concavity; see
\cite[Lemma 1.1]{Branden2014}. On the positive support, consecutive
coefficient ratios decrease, which gives the displayed inequality.
If its right side is positive, interval support makes its left factors
positive too; otherwise the inequality is immediate. This is the
defect-sequence argument of \cite[Lemma 5.3]{BercziFeher2026}.
\end{proof}

Fix $K=k-3\geq1$ and abbreviate $D_i=D(r+i,K)$. The linear identities
above express the forcing coefficient as
\begin{equation}\label{eq:expansion}
2\rho_{k,r}=\sum_{i=-5}^{1}p_i(k,r)D_i,
\end{equation}
with seven polynomials of total degree at most four. The following
recursion specifies all of them without relying on a fitted formula.
Let $F_t(c_1,\ldots,c_m;r)$ denote the formal Newton coefficient of
$\prod_{\nu=1}^m(x+c_\nu)S_{K+t}(x)$, reduced to the $D_i$ basis.
For $m=0$, start with $D(r,K+t)+D(r+1,K+t)$ and apply
\eqref{eq:defect} $t$ times. Then use
\begin{equation}\label{eq:F}
\begin{aligned}
F_t(c_1,\ldots,c_m;r)={}&(r+c_m)F_t(c_1,\ldots,c_{m-1};r)\\
&+rF_t(c_1,\ldots,c_{m-1};r-1).
\end{aligned}
\end{equation}
Substituting \eqref{eq:first} into \eqref{eq:residual} makes the right
side of \eqref{eq:expansion} exactly
\begin{equation}\label{eq:spec}
\begin{aligned}
&2F_1(2-k,0,3-k;r)-2kF_2(2-k;r)+F_1(3-k,2-k;r)\\
&\quad+2F_0(1,3-k,0,4-k;r)-2(k-1)F_1(1,3-k;r)\\
&\quad+F_0(1,4-k,3-k;r).
\end{aligned}
\end{equation}
Each defect reduction or linear multiplication raises total coefficient
degree by at most one; counting also the scalar factors in
\eqref{eq:spec} proves the degree bound four.

Put $L_{k,r}=\rho_{k,r}^2-\rho_{k,r-1}\rho_{k,r+1}$ and expand
\begin{equation}\label{eq:quadratic}
4L_{k,r}=\sum_{(a,b)\in\mathcal P}q_{a,b}(k,r)D_aD_b,
\end{equation}
where $\sort(a,b)=(\min(a,b),\max(a,b))$ and
\begin{equation}\label{eq:pairs}
\begin{aligned}
\mathcal P={}&\{\sort(a,b):-5\leq a,b\leq1\}\\
&\cup\{\sort(a-1,b+1):-5\leq a,b\leq1\}.
\end{aligned}
\end{equation}
All ranges in \eqref{eq:pairs} are integral. There are 41 pairs, including
every term of the central and adjacent products. In particular, no term
is deleted on the basis of a parameter-dependent support assumption.

For each trace $s=a+b$, order its pairs $(a_i,b_i)$ from the center
outward, that is, by increasing $b_i-a_i$. Set
\[
U_i=D_{a_i}D_{b_i},\qquad U_{t+1}=0,\qquad
W_i(k,r)=\frac14\sum_{j=1}^i q_{a_j,b_j}(k,r).
\]
Lemma~\ref{lem:defectLC} gives $U_i\geq U_{i+1}$. Abel summation on
each trace is the exact identity
\begin{equation}\label{eq:abel}
L_{k,r}=\sum_s\sum_{i=1}^t W_i(k,r)(U_i-U_{i+1}).
\end{equation}
If $U_i-U_{i+1}>0$, then $U_i>0$, so the support of $D$ implies
\[
X=r+a_i-(k-2)\geq0,\qquad
Y=2k-6-(r+b_i)\geq0.
\]
Solving these two equations gives
\begin{equation}\label{eq:substitution}
k=X+Y+b_i-a_i+4,\qquad r=2X+Y+b_i-2a_i+2.
\end{equation}

For each of the 41 prefixes, substitution \eqref{eq:substitution}
gives a nonzero polynomial in $X,Y$ all of whose nonzero coefficients
are positive. The complete inventory is Table~\ref{tab:prefixes}; the
attached certificate gives the 1,506 rational coefficients, with their
monomial exponents. The attached program reconstructs them from
\eqref{eq:defect}, \eqref{eq:F}, and \eqref{eq:spec}, forms both products
in \eqref{eq:quadratic}, and compares every coefficient exactly.

Thus each $W_i$ is nonnegative whenever its Abel difference is positive.
If that difference is zero, its summand vanishes. Equation
\eqref{eq:abel} proves \eqref{eq:main} for all $k\geq4$ and $r\geq0$.
At $k=4$, $P_1(t)=t^2$ and the same argument applies. At $r=0$ and
outside the strict support the minors are zero; at the support endpoints
they are positive squares. This completes the proof of
Theorem~\ref{thm:main}.

The certificate uses coefficient positivity on the nonnegative quadrant.
For comparison, Zeng's proof of a two-parameter spectral inequality uses exact
positive Bernstein coefficients on a compact parameter box \cite{Zeng2026Lens}
(Theorem 10, equations (8.25)--(8.27)). In both
settings, an exact polynomial identity and the sign of a basis on its
specified domain justify an infinite parameter range. Here the basis
is $X^uY^v$, and the domain is determined by permutation support.

\section{What remains for the second mode}

Write $a_r=B_{k-1,2,r}$ and $g_r=(r+1)a_r+r a_{r-1}$. The known
recurrence \eqref{eq:recurrence} says
\begin{equation}\label{eq:induction}
B_{k,2,r+1}=g_r+\rho_{k,r}.
\end{equation}
If $a$ is nonnegative and log-concave with interval support, then so is
$g$. Indeed, for $r\geq1$ its minor is
\begin{align*}
g_r^2-g_{r-1}g_{r+1}={}&
(r+1)^2a_r^2-r(r+2)a_{r-1}a_{r+1}\\
&+r(r+1)a_{r-1}a_r-(r-1)(r+2)a_{r-2}a_{r+1}\\
&+r^2a_{r-1}^2-(r^2-1)a_{r-2}a_r.
\end{align*}
By decreasing consecutive ratios and the zero-boundary convention,
the three lines are at least $a_r^2$, $2a_{r-1}a_r$, and
$a_{r-1}^2$, respectively. The case $r=0$ is immediate, and
\eqref{eq:linear} gives interval support.

Consequently one sufficient remaining assertion is
\begin{equation}\label{eq:mixed}
2g_r\rho_{k,r}-g_{r-1}\rho_{k,r+1}-g_{r+1}\rho_{k,r-1}\geq0
\qquad(k\geq5,\ r\geq0).
\end{equation}
This assertion is unproved here. Adding its left side to the separate
minors of $g$ and $\rho$ gives the minor of $g+\rho$. The base case
$k=4$ follows from $A_{3,2}=0$, $R_4=\Delta A_{4,2}$, and
$B_{4,2,0}=0$. For the induction, the positive supports of $a$ and $g$
are $[k-2,2k-2]$ and $[k-2,2k-1]$, respectively, so they also have the
required overlap with the forcing support. Thus \eqref{eq:mixed}
would suffice for the second mode; no equivalence with the original
conjecture is asserted.

The numerical extent of the supplementary checks is explicit. A separate
implementation rebuilt \eqref{eq:chern} for $0\leq d\leq90$ and checked
3,690 residual Newton entries and their minors for $4\leq k\leq44$,
along with 3,690 recurrence evaluations. It also found no negative
mixed expression in the 2,000 cases $5\leq k\leq44$, $0\leq r\leq2k$.
These are exact finite observations, not a proof of \eqref{eq:mixed}.
For example, the first residual row is
\[
(\rho_{4,r})_{r=0}^{8}=(0,0,27,212,660,1000,735,210,0).
\]

\appendix
\section{Certificate data and reproducibility}\label{app:certificate}

Two supplementary files are embedded in this PDF: \texttt{certificate.json}
contains all seven expansion polynomials and all 41 complete prefixes;
\texttt{certificate.py} reconstructs them using only the Python 3 standard
library. They can be saved from a PDF reader's attachment panel, or with
\texttt{pdfdetach -saveall main.pdf}. In the directory containing both files,
run
\begin{center}
\verb|python3 certificate.py --check certificate.json|
\end{center}
The program checks the complete pair set, degree bounds, coefficient
identities, and all signs. It reports seven expansion polynomials,
41 prefixes, and 1,506 positive coefficients. The variable names in each
JSON field specify whether the exponent pairs refer to $(k,r)$ or
$(X,Y)$. Coefficients are exact rational strings; omitted monomials have
coefficient zero. The attached program can also reconstruct the entire
JSON file with \verb|--write certificate.json|.

The initial symbolic generation used SymPy \cite{Meurer2017}. Two separate
standard-library implementations checked the finite algebra: one used
exact evaluation on a $9\times9$ tensor grid for every prefix, and the
other compared full sparse polynomial dictionaries. The grid calculation
has a proof of completeness: each actual prefix has total degree at most
eight by \eqref{eq:expansion}--\eqref{eq:quadratic}; affine substitution
does not increase that bound. The reported polynomials satisfy the same
bound. A polynomial separately of degree at most eight in each variable
that vanishes at all $(X,Y)\in\{0,\ldots,8\}^2$ is identically zero,
by applying univariate interpolation twice. Hence the 3,321 exact
comparisons prove identities rather than extrapolating finite inequalities.
The polynomial certificate is the infinite-range justification; the
direct-product checks reported above are additional consistency tests.

\begin{table}[htbp]
\centering
\caption{Complete prefix inventory, read down the left block and then the right.
Each row describes $W_{a,b}(X,Y)$ after \eqref{eq:substitution}; $d$ is total degree,
$m$ counts nonzero coefficients, and $c_{\min}$ is their minimum. The attached
JSON contains all 1,506 coefficients.}\label{tab:prefixes}
\small
\begin{tabular}{rrrrr@{\qquad}rrrrr}
\toprule
Trace & Pair & $d$ & $m$ & $c_{\min}$ & Trace & Pair & $d$ & $m$ & $c_{\min}$\\
\midrule
$-10$ & $(-5,-5)$ & 8 & 42 & $1$ & $-4$ & $(-5,1)$ & 8 & 40 & $1$ \\
$-10$ & $(-6,-4)$ & 6 & 28 & $1$ & $-4$ & $(-6,2)$ & 6 & 28 & $210$ \\
$-9$ & $(-5,-4)$ & 8 & 42 & $5$ & $-3$ & $(-2,-1)$ & 8 & 42 & $40$ \\
$-9$ & $(-6,-3)$ & 6 & 28 & $10$ & $-3$ & $(-3,0)$ & 8 & 42 & $10$ \\
$-8$ & $(-4,-4)$ & 8 & 42 & $15$ & $-3$ & $(-4,1)$ & 8 & 40 & $5$ \\
$-8$ & $(-5,-3)$ & 8 & 42 & $10$ & $-3$ & $(-5,2)$ & 6 & 28 & $120$ \\
$-8$ & $(-6,-2)$ & 6 & 28 & $45$ & $-2$ & $(-1,-1)$ & 8 & 42 & $\tfrac{43}{4}$ \\
$-7$ & $(-4,-3)$ & 8 & 42 & $40$ & $-2$ & $(-2,0)$ & 8 & 42 & $10$ \\
$-7$ & $(-5,-2)$ & 8 & 42 & $10$ & $-2$ & $(-3,1)$ & 8 & 40 & $10$ \\
$-7$ & $(-6,-1)$ & 6 & 28 & $120$ & $-2$ & $(-4,2)$ & 6 & 28 & $45$ \\
$-6$ & $(-3,-3)$ & 8 & 42 & $50$ & $-1$ & $(-1,0)$ & 8 & 42 & $5$ \\
$-6$ & $(-4,-2)$ & 8 & 42 & $45$ & $-1$ & $(-2,1)$ & 8 & 40 & $10$ \\
$-6$ & $(-5,-1)$ & 8 & 42 & $5$ & $-1$ & $(-3,2)$ & 6 & 28 & $10$ \\
$-6$ & $(-6,0)$ & 6 & 28 & $210$ & $0$ & $(0,0)$ & 8 & 42 & $1$ \\
$-5$ & $(-3,-2)$ & 8 & 42 & $75$ & $0$ & $(-1,1)$ & 8 & 40 & $1$ \\
$-5$ & $(-4,-1)$ & 8 & 42 & $24$ & $0$ & $(-2,2)$ & 6 & 28 & $1$ \\
$-5$ & $(-5,0)$ & 8 & 42 & $1$ & $1$ & $(0,1)$ & 8 & 39 & $1$ \\
$-5$ & $(-6,1)$ & 6 & 28 & $252$ & $1$ & $(-1,2)$ & 6 & 27 & $2$ \\
$-4$ & $(-2,-2)$ & 8 & 42 & $\tfrac{77}{2}$ & $2$ & $(1,1)$ & 8 & 25 & $1$ \\
$-4$ & $(-3,-1)$ & 8 & 42 & $45$ & $2$ & $(0,2)$ & 6 & 25 & $2$ \\
$-4$ & $(-4,0)$ & 8 & 42 & $5$ &  & & & &  \\
\bottomrule
\end{tabular}
\end{table}

\paragraph{Computational and AI assistance.}
AI systems assisted with proof exploration, program construction, source
screening, and manuscript drafting. A separate AI review pass reconstructed
the mathematical argument and verified the certificate with a separate
exact-arithmetic implementation. These are reproducible computational
checks and a separate review pass; no proof-assistant formalization or
journal peer review is claimed.

% The audited build attaches certificate.json and certificate.py using pdfattach.
\bibliographystyle{abbrvnat}
\bibliography{references}
\end{document}
